% The sequence type hierarchy (figure in sequences.tex).  Abstract
% types in plain face, instantiable ones in bold; arrows point at
% supertypes.
\begin{tikzpicture}[
  x=2.2cm, y=1.2cm,
  every node/.style={ellipse, draw, inner xsep=6pt, inner ysep=2pt, font=\small},
  inst/.style={font=\small\bfseries},
  sup/.style={-{Stealth[length=6pt]}, thick}]
  \node (sequence)  at (2.2,3)   {sequence};
  \node (vector)    at (0.9,2)   {vector-type};
  \node (list)      at (3.3,2)   {list-type};
  \node (svector)   at (0.7,1) [inst] {simple-vector};
  \node (pair)      at (2.7,1) {pair};
  \node (null)      at (4.3,1) {null-type};
  \node (string)    at (0.5,0)   [inst] {string};
  \node (cons)      at (1.9,0) [inst] {cons-pair};
  \node (lazy)      at (3.9,0) [inst] {lazy-cons-pair};
  \draw[sup] (vector)  -- (sequence);
  \draw[sup] (list)    -- (sequence);
  \draw[sup] (svector) -- (vector);
  \draw[sup] (pair)    -- (list);
  \draw[sup] (null)    -- (list);
  \draw[sup] (string)  -- (svector);
  \draw[sup] (cons)    -- (pair);
  \draw[sup] (lazy)    -- (pair);
\end{tikzpicture}
